\documentclass[11pt]{article}

\usepackage[a4paper,margin=2.5cm]{geometry}
\usepackage[utf8]{inputenc}
\usepackage[T1]{fontenc}
\usepackage[french]{babel}
\usepackage{amsmath,amssymb}
\usepackage{booktabs}
\usepackage{array}
\usepackage{hyperref}

\title{Dimensionnement PV + Batterie + Réseau\\
\large Modélisation mathématique (MILP)}
\author{}
\date{}

\begin{document}
\maketitle

\section{Contexte}

On cherche à dimensionner, pour une maison individuelle, un système composé de panneaux
photovoltaïques et d'une batterie stationnaire, connectés au réseau électrique, de façon à
\textbf{minimiser le coût total annualisé} tout en \textbf{couvrant la demande électrique du
foyer à chaque heure de l'année}. Le problème est résolu sur un horizon d'un an au pas
horaire (8760 heures), sous la forme d'un programme linéaire en nombres entiers mixtes
(MILP) : deux variables de dimensionnement sont entières (nombre de panneaux, nombre de
modules de batterie), le reste du problème est linéaire continu.

\section{Notations}

\subsection{Ensemble}

\[
t \in \mathcal{T} = \{1, \dots, 8760\} \quad \text{(les heures de l'année)}
\]

\subsection{Paramètres}

\begin{center}
\begin{tabular}{@{}lll@{}}
\toprule
Symbole & Description & Unité \\
\midrule
$D_t$ & Demande électrique du foyer à l'heure $t$ & kWh \\
$\pi_t$ & Production PV d'\emph{un seul} panneau à l'heure $t$ & kWh \\
$c^{ach}_t$ & Prix d'achat de l'électricité à l'heure $t$ & €/kWh \\
$c^{vte}_t$ & Prix de vente du surplus à l'heure $t$ (constant) & €/kWh \\
$c_{abo}$ & Abonnement annuel au réseau & € \\
$c_{pv}$ & Coût d'un panneau installé & € \\
$c^{kW}_{b}$, $c^{kWh}_{b}$ & Coût unitaire de puissance / capacité batterie & €/kW, €/kWh \\
$\eta^{ch}$, $\eta^{dc}$ & Rendements de charge / décharge de la batterie & -- \\
$e_{mod}$, $p_{mod}$ & Capacité / puissance d'un module de batterie standard & kWh, kW \\
$r$ & Taux d'actualisation & -- \\
$n^{pv}_{vie}$, $n^{b}_{vie}$ & Durée de vie (PV, batterie) & années \\
$CRF_{pv}$, $CRF_{b}$ & Facteurs d'annualisation du capex (PV, batterie) & -- \\
$\tau$ & Plafond imposé sur la part de la demande couverte par le réseau & -- \\
\bottomrule
\end{tabular}
\end{center}

Le facteur d'annualisation du capital (\emph{Capital Recovery Factor}) répartit un
investissement payé une seule fois sur sa durée de vie, à un taux d'actualisation $r$ :
\begin{equation}
CRF(r, n) = \frac{r(1+r)^n}{(1+r)^n - 1}
\qquad CRF_{pv} = CRF(r, n^{pv}_{vie}), \quad CRF_{b} = CRF(r, n^{b}_{vie})
\end{equation}

\subsection{Variables de décision}

\paragraph{Dimensionnement (entiers).}
\begin{align*}
n_{pv} &\in \mathbb{N} && \text{nombre de panneaux PV installés} \\
n_{b} &\in \mathbb{N} && \text{nombre de modules de batterie installés}
\end{align*}

\paragraph{Dimensionnement dérivé (continu, lié aux entiers ci-dessus).}
\begin{align*}
P^{max}_{b} &\geq 0 && \text{puissance totale de la batterie (kW)} \\
E^{max}_{b} &\geq 0 && \text{capacité totale de la batterie (kWh)}
\end{align*}

\paragraph{Flux horaires (continus, $\geq 0$, en kWh/h), pour tout $t \in \mathcal{T}$.}

\begin{center}
\begin{tabular}{@{}lll@{}}
\toprule
Symbole & Nom dans le code & Description \\
\midrule
$a_t$ & \texttt{achat\_electricite} & Électricité achetée au réseau \\
$v_t$ & \texttt{vente\_electricite} & Électricité vendue au réseau \\
$f^{pv \to h}_t$ & \texttt{pv\_maison} & PV $\to$ maison (autoconsommation directe) \\
$f^{pv \to b}_t$ & \texttt{pv\_batterie} & PV $\to$ batterie \\
$f^{pv \to g}_t$ & \texttt{pv\_grid} & PV $\to$ réseau (vente de surplus) \\
$f^{g \to b}_t$ & \texttt{grid\_batterie} & Réseau $\to$ batterie \\
$f^{g \to h}_t$ & \texttt{grid\_maison} & Réseau $\to$ maison \\
$f^{b \to h}_t$ & \texttt{batterie\_maison} & Batterie $\to$ maison \\
$f^{b \to g}_t$ & \texttt{batterie\_grid} & Batterie $\to$ réseau \\
$ch_t$ & \texttt{echarge} & Charge totale de la batterie \\
$dc_t$ & \texttt{edecharge} & Décharge totale de la batterie \\
$E_t$ & \texttt{estock\_batterie} & État de charge de la batterie \\
\bottomrule
\end{tabular}
\end{center}

\section{Fonction objectif}

Minimiser le coût total annuel = abonnement + capex annualisé (PV et batterie) + coût net
d'électricité :
\begin{equation}
\min \; c_{abo} \cdot \delta
\;+\; n_{pv}\, c_{pv}\, CRF_{pv}
\;+\; \big(P^{max}_{b}\, c^{kW}_{b} + E^{max}_{b}\, c^{kWh}_{b}\big)\, CRF_{b}
\;+\; \sum_{t \in \mathcal{T}} \Big( a_t \cdot c^{ach}_t - v_t \cdot c^{vte}_t \Big)
\end{equation}
où $\delta \in \{0,1\}$ vaut $1$ si le foyer reste connecté au réseau (option, voir
section~\ref{sec:extensions}) ; par défaut $\delta = 1$ (fixe, connexion obligatoire).

\section{Contraintes}

\paragraph{Bilan de production PV.} La production totale (nombre de panneaux $\times$
production unitaire) se répartit entre autoconsommation directe, charge de la batterie et
injection réseau :
\begin{equation}
\pi_t \cdot n_{pv} = f^{pv\to h}_t + f^{pv\to b}_t + f^{pv\to g}_t \qquad \forall t
\end{equation}

\paragraph{Bilan des achats réseau.} L'électricité achetée sert à la maison et/ou à
charger la batterie (le second terme disparaît si l'arbitrage batterie est désactivé, voir
section~\ref{sec:extensions}) :
\begin{equation}
a_t = f^{g\to b}_t + f^{g\to h}_t \qquad \forall t
\end{equation}

\paragraph{Bilan des ventes réseau.} L'électricité vendue provient du PV et/ou de la
batterie :
\begin{equation}
v_t = f^{b\to g}_t + f^{pv\to g}_t \qquad \forall t
\end{equation}

\paragraph{Bilan de charge/décharge de la batterie.}
\begin{align}
ch_t &= f^{g\to b}_t + f^{pv\to b}_t \qquad \forall t \\
dc_t &= f^{b\to h}_t + f^{b\to g}_t \qquad \forall t
\end{align}

\paragraph{Dynamique de l'état de charge (cyclique).} Avec pertes de rendement, et
bouclage de l'année sur elle-même (l'état de fin d'année alimente le début, plutôt qu'une
valeur initiale figée arbitrairement à 0) :
\begin{equation}
E_t = E_{t-1} + ch_t \cdot \eta^{ch} - \frac{dc_t}{\eta^{dc}}
\qquad \forall t, \quad \text{avec } E_0 \equiv E_{8760}
\end{equation}

\paragraph{Bornes de capacité et de puissance.}
\begin{align}
0 &\leq E_t \leq E^{max}_{b} \qquad \forall t \\
ch_t &\leq P^{max}_{b} \qquad dc_t \leq P^{max}_{b} \qquad \forall t
\end{align}

\paragraph{Discrétisation de la batterie.} Capacité et puissance sont des multiples
entiers d'un module standard, plutôt que des variables continues (réalisme commercial) :
\begin{equation}
E^{max}_{b} = n_{b} \cdot e_{mod} \qquad P^{max}_{b} = n_{b} \cdot p_{mod}
\end{equation}

\paragraph{Couverture de la demande.} Contrainte centrale du problème : la demande est
satisfaite à chaque heure par le PV direct, le réseau ou la batterie :
\begin{equation}
D_t = f^{pv\to h}_t + f^{g\to h}_t + f^{b\to h}_t \qquad \forall t
\end{equation}

\paragraph{Plafond d'achat réseau (scénario de transition énergétique).} Optionnel : borne
la part de la demande annuelle qui peut être couverte par des achats au réseau, maison et
charge de la batterie confondues (peu importe la destination finale de l'énergie achetée) :
\begin{equation}
\sum_{t \in \mathcal{T}} a_t \;\leq\; \tau \cdot \sum_{t \in \mathcal{T}} D_t
\end{equation}
$\tau = 1$ correspond au marché libre (aucune contrainte), $\tau = 0$ à une autonomie
totale imposée (aucun achat réseau, à aucune heure de l'année). Cette formulation ne peut
pas être contournée par un chargement de la batterie depuis le réseau, contrairement à une
contrainte qui ne porterait que sur l'autoconsommation directe.

\subsection{Extensions optionnelles}
\label{sec:extensions}

\paragraph{Restriction anti-arbitrage (activée par défaut).} Pour éviter que le modèle
n'exploite un écart entre un prix d'achat variable (marché spot, pouvant devenir négatif)
et un prix de vente fixe garanti -- une stratégie plus proche du trading financier que de
l'autoconsommation -- la charge de la batterie depuis le réseau peut être interdite :
\begin{equation}
f^{g\to b}_t = 0 \qquad \forall t
\end{equation}
Dans ce cas, le bilan de charge de la batterie se réduit à $ch_t = f^{pv\to b}_t$ : la
batterie ne stocke que du surplus PV, revendable au réseau plus tard sans que ce soit du
trading.

\paragraph{Déconnexion du réseau (optionnelle).} Un binaire $\delta \in \{0,1\}$ peut
rendre l'abonnement réseau optionnel :
\begin{equation}
a_t \leq M \cdot \delta \qquad v_t \leq M \cdot \delta \qquad \forall t
\end{equation}
où $M$ est une borne supérieure large sur les flux horaires réseau. Si $\delta = 0$, le
foyer doit être strictement autosuffisant à chaque heure (pas seulement en moyenne
annuelle) et ne paie pas l'abonnement.

\section{Résumé des hypothèses clés}

\begin{itemize}
\item Année météo et profil de demande \textbf{représentatifs}, pas des mesures réelles
du foyer.
\item Le capex est \textbf{annualisé} via $CRF$ (taux d'actualisation $r$, durées de vie
$n^{pv}_{vie}$, $n^{b}_{vie}$) plutôt que comparé tel quel à des économies sur un an.
\item La batterie est \textbf{discrétisée} par modules entiers de taille fixe
($e_{mod}$, $p_{mod}$), pas dimensionnée en continu.
\item Le plafond $\tau$ porte sur le \textbf{total annuel} des achats réseau, pas sur
chaque heure individuellement.
\item Aucune borne de surface de toit ni de budget maximal n'est imposée sur $n_{pv}$ ou
$n_{b}$.
\item Les prix d'investissement (PV, batterie), le taux d'actualisation et les durées de
vie sont incertains -- une analyse de sensibilité (hors modélisation, voir le notebook)
teste la robustesse des conclusions à ces hypothèses.
\end{itemize}

\end{document}
